% ================= 第 4 章 (前半): 4.1 + 4.2 =================

\section{超收敛伴随形式}\label{sec:4}

我们证明: 当 $f$ 是 $p$ 处序形式时 (因而特别是如 \ref{sec:3.3} 时), 表示 $\sigma_p(f)$ 分裂当且仅当与 $f$ 相关联的某个斜率 $k-1$ 的模形式具有 $\theta^{k-1}(g)$ 的形式, 其中 $g$ 是权 $2-k$ 的序超收敛形式 (引言的定理 \ref{thm:1.1.3}).

\subsection{\mbox{}}\label{sec:4.1}

我们简要回顾"有限斜率"和"$U_p$ 的特征值"这两个经典概念的表示论解释.
假设 $f$ 是如 \ref{sec:3.3} 的新形式, 因而 $\pi_p(f)$ 是主级数.
用 \ref{sec:3.3} 的记号, 回顾:
\[
\pi_p(f) = \lcind_{B(\Qp)}^{\GL_2(\Qp)} \chi_1z^{k-2} \otimes \chi_2.
\]
新形式 $f$ 是有限斜率的当且仅当特征标 $\chi_1z^{k-2}$ 或 $\chi_2$ 至少有一个是非分歧的.
若两者都非分歧, 则 $f$ 的导子与 $p$ 互素, 且是 $T_p$ 的特征向量, 其 $p$ 处 Hecke 多项式为 $X^{2} - \bigl(\chi_1(p)p^{k-1} + \chi_2(p)\bigr)X + \chi(p)p^{k-1}$ (利用 $\chi_1\chi_2 = \chi_p^{-1}$ 和 $\chi_p(p)^{-1} = \chi(p)$).
此时形式 $f$ 给出两个 $p$ 处"$p$-稳定化"的旧形式, 有限斜率, 且分别是 $U_p$ 关于特征值 $\alpha_p := \chi_2(p)$ 和 $\beta_p := p^{k-1}\chi_1(p)$ 的特征向量.
例如, $f(z) - \beta_p f(pz)$ 是特征值 $\alpha_p$ (且斜率 $k-1$) 的旧形式.
若特征标 $\chi_1z^{k-2}$ 或 $\chi_2$ 只有一个非分歧, 必要时应用如 \ref{sec:3.3} 的缠结, 可以假设 $\chi_1z^{k-2}$ 非分歧.
在这种情况下, $p$ 整除 $f$ 的导子, 且 $f$ 是 $U_p$ 的特征向量, 特征值为 $p^{k-1}\chi_1(p)$.

设 $\chi_0$ 为 $\chi$ 的 $p$-部分 (它是 $(\mathbb{Z}/p^{r}\mathbb{Z})^{\times}$ 的特征标), 记 $\tilde{f}$ 为与扭曲模形式 $f \otimes \chi_0^{-1}$ 相关联的新形式 (它不是新的, 参见下文).
则有:
\[
\pi_p(\tilde{f}) = \lcind_{B(\Qp)}^{\GL_2(\Qp)} \chi_1z^{k-2}\chi_0 \otimes \chi_2\chi_0 = \lcind_{B(\Qp)}^{\GL_2(\Qp)} \chi_2|\,|\chi_0 \otimes \chi_1z^{k-2}|\,|^{-1}\chi_0
\]
(其中已把 $\chi_0$ 延拓为 $\Qp^{\times}$ 的特征标, 令 $\chi_0(p) = 1$).
特征标 $\chi_2|\,|\chi_0$ 非分歧, 故前述讨论表明 $\tilde{f}$ 对 $U_p$ 的特征值 $\chi_2(p)$ 是有限斜率 ($= k-1$) 的 (此时显然 $\tilde{f} \neq f \otimes \chi_0^{-1}$, 因为后一形式被 $U_p$ 零化).
在余下部分, 我们将使用 $\tilde{f}$ 的众所周知公式:
\[
\tilde{f}(z) = (f\,|\,w_{p^{r}})(z) := (Ncz + p^{r}d)^{k} f\!\left(\frac{p^{r}az + b}{Ncz + p^{r}d}\right),
\]
其中 $w_{p^{r}} := \begin{pmatrix} p^{r}a & b \\ Nc & p^{r}d \end{pmatrix}$, $a,b,c,d \in \mathbb{Z}$, $b \equiv 1 \bmod p^{r}$, $a \equiv 1 \bmod M$, 且 $p^{r}ab - Mcd = 1$ (此公式可用 adele 计算验证, 也可以用例如 \cite[\S6]{ref:30} 的显式公式).

\subsection{\mbox{}}\label{sec:4.2}

我们给出模曲线上的 Gauss--Manin 联络和算子 $\theta^{k-1}$ 的一些预备知识 (参见例如 \cite{ref:13} 和 \cite{ref:12} 的开头).
在下面的所有内容中, 我们不假说明地使用如下事实: $\Qp$ 上紧概形 $X$ 上凝聚层的上同调等于从 $X$ 得到的刚性解析簇上"逆像"层的上同调 (刚性 GAGA).

先固定整数 $N > 4$, 记 $\pi : E \to X_1(N)$ 为 $\Qp$ 上紧的通常模曲线 $X_1(N)$ (与同余子群 $\Gamma_1(N)$ 相关联) 之上的泛广义椭圆曲线, $C \subset X_1(N)$ 为尖点闭子概形, $\widetilde{C} := \pi^{-1}(C)$, $\omega := \pi_{*}\Omega^{1}_{E/X_1(N)}(\log \widetilde{C})$, 以及:
\[
H_{k-2} := \Sym^{k-2}R^{1}\pi_{*}\Omega^{\bullet}_{E/X_1(N)}(\log \widetilde{C})
\]
(其中 $H_0 = \mathcal{O}_{X_1(N)}$).
回顾: 层 $\omega$ 在 $X_1(N)$ 上可逆, 层 $R^{1}\pi_{*}\Omega^{\bullet}_{E/X_1(N)}(\log \widetilde{C})$ 秩 $2$ 局部自由, 层 $H_{k-2}$ 秩 $k-1$ 局部自由.
$H_{k-2}$ 上还有由局部直和因子子 $\mathcal{O}_{X_1(N)}$-模 $\mathrm{Fil}^{i}H_{k-2}$ 组成的规范递降滤过 $(\mathrm{Fil}^{i}H_{k-2})_{i \in \mathbb{Z}}$ (Hodge 滤过), 使得:
\[
\mathrm{Fil}^{i}H_{k-2} = H_{k-2},\quad i \le 0;
\qquad
\mathrm{Fil}^{i}H_{k-2}\ \text{秩 } k-1-i \text{ 局部自由},\quad 1 \le i \le k-2;
\qquad
\mathrm{Fil}^{i}H_{k-2} = 0,\quad i \ge k-1.
\]
此滤过由 $\mathrm{Fil}^{1}H_1 = \omega$ 以及 $\mathrm{Fil}^{i}H_r \cdot \mathrm{Fil}^{j}H_s = \mathrm{Fil}^{i+j}H_{r+s}$ ($0 \le i \le r$, $0 \le j \le s$) 决定.
因此 $\mathrm{Fil}^{k-2}H_{k-2} = \omega^{k-2}$, 且更一般地可以证明: 对 $0 \le i \le k-2$, $\mathrm{gr}^{i}H_{k-2} \simeq \omega^{2i-k+2}$.
Gauss--Manin 联络延拓为对数 de Rham 复形:
\[
\nabla : H_{k-2} \to H_{k-2} \otimes_{\mathcal{O}_{X_1(N)}} \Omega^{1}_{X_1(N)/\Qp}(\log C),
\]
并满足 Griffiths 横截性, 即对 $i \in \mathbb{Z}$ 定义子复形:
\[
\nabla : \mathrm{Fil}^{i}H_{k-2} \to \mathrm{Fil}^{i-1}H_{k-2} \otimes_{\mathcal{O}_{X_1(N)}} \Omega^{1}_{X_1(N)/\Qp}(\log C).
\]
此外, 它对 $0 \le i \le k-2$ 的梯度诱导同构:
\begin{equation}
\label{eq:17}
\nabla : \omega^{2i-k+2} \xrightarrow{\sim} \omega^{2i-k} \otimes_{\mathcal{O}_{X_1(N)}} \Omega^{1}_{X_1(N)/\Qp}(\log C)
\end{equation}
(若 $k = 2i$, 它就是 Kodaira--Spencer 同构 $\omega^{2} \simeq \Omega^{1}_{X_1(N)/\Qp}(\log C)$).
若记 $\mathrm{Fil}^{i}H^{1}\bigl(X_1(N),H_{k-2} \otimes \Omega^{\bullet}(\log)\bigr)$ 为下列上同调群:
\[
H^{1}\bigl(X_1(N),\ \mathrm{Fil}^{i}H_{k-2} \to \mathrm{Fil}^{i-1}H_{k-2} \otimes_{\mathcal{O}_{X_1(N)}} \Omega^{1}_{X_1(N)/\Qp}(\log C)\bigr),
\]
由 \eqref{eq:17} 经显然的上同调计算可验证: 有注入 $\mathrm{Fil}^{1}H^{1}\bigl(X_1(N),H_{k-2} \otimes \Omega^{\bullet}(\log)\bigr) \hookrightarrow H^{1}\bigl(X_1(N),H_{k-2} \otimes \Omega^{\bullet}(\log)\bigr)$, $\mathrm{Fil}^{i}H^{1}\bigl(X_1(N),H_{k-2} \otimes \Omega^{\bullet}(\log)\bigr)$ 在 $1 \le i \le k-1$ 时为常数, 且:
\[
\mathrm{Fil}^{k-1}H^{1}\bigl(X_1(N),H_{k-2} \otimes \Omega^{\bullet}(\log)\bigr)
\simeq H^{0}\bigl(X_1(N),\omega^{k-2} \otimes_{\mathcal{O}_{X_1(N)}} \Omega^{1}_{X_1(N)/\Qp}(\log C)\bigr)
\simeq H^{0}(X_1(N),\omega^{k})
\]
(经 Kodaira--Spencer).
最后回顾复合映射:
\[
\psi : \mathrm{Fil}^{1}H_{k-2} \hookrightarrow H_{k-2} \xrightarrow{\nabla} H_{k-2} \otimes \Omega^{1}(\log C) \twoheadrightarrow \bigl(H_{k-2}/\mathrm{Fil}^{k-2}H_{k-2}\bigr) \otimes \Omega^{1}(\log C)
\]
是同构 (利用 \eqref{eq:17}), 它诱导 $X_1(N)$ 上层的一个规范 $\Qp$-线性态射:
\[
\omega^{-k+2} \simeq \mathrm{gr}^{0}H_{k-2} \to \mathrm{Fil}^{k-2}H_{k-2} \otimes \Omega^{1}_{X_1(N)/\Qp}(\log C) \simeq \omega^{k},
\qquad
s \mapsto \nabla \overline{s} - \psi^{-1}\bigl(\overline{\nabla(s)}\bigr),
\]
其中 $\overline{s}$ 是 $s$ 在 $H_{k-2}$ 中的任意一个提升, $\overline{\nabla(s)}$ 是 $\nabla(\overline{s})$ 在 $\bigl(H_{k-2}/\mathrm{Fil}^{k-2}H_{k-2}\bigr) \otimes \Omega^{1}_{X_1(N)/\Qp}(\log C)$ 中的像.
若 $U \subset X_1(N)$ 是与 $X_1(N)$ 相关联的刚性簇的任意连通开集且包含无穷远尖点, 则诱导的映射 $H^{0}(U,\omega^{2-k}) \to H^{0}(U,\omega^{k})$ 在 $q$-展开上等同于 $\frac{(-1)^{k}}{k!}\theta^{k-1}$, 其中 $\theta$ 是 $\Qp[[q]]$ 上的算子 $q(d/dq)$ (细节参见 \cite[\S9]{ref:12}).
为记号简便, 对 $X_1(N)$ 的任何刚性开集 $U \subset X_1(N)$ (不必连通), 仍记 $\theta^{k-1} : H^{0}(U,\omega^{2-k}) \to H^{0}(U,\omega^{k})$ 为前面的映射乘以 $k!(-1)^{k}$ 所得.

\begin{lemma}\label{lem:4.2.1}
设 $U$ 是 $X_1(N)$ (看作 $\Qp$ 上刚性簇) 的拟 Stein 开集.
注入:
\[
H^{0}(U,\omega^{k}) \simeq H^{0}(U,\omega^{k-2} \otimes \Omega^{1}(\log C)) \hookrightarrow H^{0}(U,H_{k-2} \otimes \Omega^{1}(\log C))
\]
诱导同构:
\[
\frac{H^{0}(U,\omega^{k})}{\theta^{k-1}H^{0}(U,\omega^{2-k})} \xrightarrow{\sim} \frac{H^{0}(U,H_{k-2} \otimes \Omega^{1}(\log C))}{\nabla H^{0}(U,H_{k-2})} \simeq H^{1}\bigl(U,H_{k-2} \otimes \Omega^{\bullet}(\log)\bigr).
\]
\end{lemma}

\begin{proof}
第二个同构由如下事实得出: 对 $U$ 上任何凝聚层 $F$ 和一切 $q > 0$, $H^{q}(U,F) = 0$.
在 $X_1(N)$ 上有层正合列 (用上述记号):
\[
0 \to \mathrm{gr}^{0}H_{k-2} \to \frac{H_{k-2}}{\mathrm{Fil}^{k-2}H_{k-2}} \to \mathrm{Fil}^{1}H_{k-2} \to 0,
\]
其中第一个映射是 $s \mapsto \overline{s} - \psi^{-1}\bigl(\overline{\nabla(s)}\bigr)$, 第二个是 $t \mapsto \psi^{-1}\bigl(\overline{\nabla(t)}\bigr)$,
它嵌进层复形的一个正合列:
\[
\begin{CD}
0 @>>> \mathrm{gr}^{0}H_{k-2} @>>> \dfrac{H_{k-2}}{\mathrm{Fil}^{k-2}H_{k-2}} @>>> \mathrm{Fil}^{1}H_{k-2} @>>> 0\\
@. @VV{\frac{(-1)^{k}}{k!}\theta^{k-1}}V @VV{\nabla}V @VV{\psi}V\\
0 @>>> \mathrm{Fil}^{k-2}H_{k-2} \otimes \Omega^{1}(\log C) @>>> H_{k-2} \otimes \Omega^{1}(\log C) @>>> \dfrac{H_{k-2}}{\mathrm{Fil}^{k-2}H_{k-2}} \otimes \Omega^{1}(\log C) @>>> 0.
\end{CD}
\]
左边复形的 $H^{1}$ 因 $U$ 拟 Stein 而等同于 $\frac{H^{0}(U,\omega^{k})}{\theta^{k-1}H^{0}(U,\omega^{2-k})}$.
相关联的长正合列则给出结果, 因为右边复形的上同调为零.
\end{proof}

现在设 $N$ 为满足 $1 \le N \le 4$ 的整数.
前面的讨论不能直接适用, 因为曲线 $X_1(N)$ 不再是精细模空间.
在这种情况下我们遵循 \cite{ref:10} 附录中描述的方法.
选择一个素数 $q > 2$, $q \nmid Np$, 并考虑与同余子群 $\Gamma_1(N) \cap \Gamma(q)$ 相关联的模曲线 $X_1(N;q)$.
$q$ 的选择使得该曲线现在是参数化"带精确 $N$ 阶点的椭圆曲线 $E$ 及同构 $\iota : E[q] \xrightarrow{\sim} (\mathbb{Z}/q\mathbb{Z})^{2}$"的模问题的精细模空间.
$\GL_2(\mathbb{Z}/q\mathbb{Z})$ 在同构 $\iota$ 的集合上的作用诱导 $\GL_2(\mathbb{Z}/q\mathbb{Z})$ 在 $X_1(N;q)$ 上的作用, 其商是 $X_1(N)$.
由于 $X_1(N;q)$ 是精细模空间, 前面的讨论 mutatis mutandis 适用于 $X_1(N;q)$.
此外, 所涉及的所有层与联络都自然 $\GL_2(\mathbb{Z}/q\mathbb{Z})$-等变.
特别地, 若在引理 \ref{lem:4.2.1} 对 $X_1(N;q)$ 的类似陈述中假设拟 Stein 开集 $U$ 是 $\GL_2(\mathbb{Z}/q\mathbb{Z})$-不变的, 则该引理的同构也是 $\GL_2(\mathbb{Z}/q\mathbb{Z})$-等变的.
